% ============================================================
%  SEKCJA 7: PROPOZYCJA MODELU OCENY DOJRZAŁOŚCI FUNKCJONALNEJ
% ============================================================

\section{Propozycja modelu oceny dojrzałości funkcjonalnej}

Rozważania prowadzone w~poprzednich sekcjach doprowadziły
do~sformułowania tezy, że~dojrzałość miodu jest zjawiskiem
ciągłym i~wielowymiarowym, i~że żaden pojedynczy parametr,
ani~morfologiczny (zasklepienie), ani~mikrobiologiczny (CFU/g),
nie~jest wystarczający do~jej wiarygodnej oceny. 
Skoro dojrzałość/dojrzewanie miodu jest zjawiskiem ciągłym winno się posługiwać wyłącznie pojęciem “stopień dojrzałości”. Wprowadzenie jakiejkolwiek klasyfikacji w tym zakresie jest nie tyle niezwykle trudne, co praktycznie niemożliwe z uwagi na naturalną zmienność miodu wynikającą co najmniej z uwarunkowań klimatycznych, pogodowych, metod prowadzenia gospodarki pasiecznej, zróżnicowania pożytków. Ograniczyć można jednak istotne ryzyko jakie dla bezpieczeństwa produktu  i jego wartości handlowej wnosi dojrzałość miodu czyli ryzyko wystąpienia fermentacji. Ta, jak powyżej  wskazano, jest zależna bezpośrenio od aktywności wody.

% W~niniejszej
% sekcji proponujemy operacyjny \textbf{model oceny dojrzałości
% funkcjonalnej miodu} (MODF), oparty na~hierarchicznym podejściu
% trójpoziomowym: od~szybkich, tanich parametrów przesiewowych
% do~zaawansowanych narzędzi metabolomicznych stosowanych
% w~przypadkach spornych. Proponowany model jest kompatybilny
% z~obowiązującymi regulacjami (Codex STAN 12-1981,
% Dyrektywa UE 2001/110/WE) i~nie wymaga ich formalnej rewizji, jedynie uzupełnienia i~reinterpretacji w~świetle aktualnej
% wiedzy naukowej.

% %,----------------------------------------------------------
% \subsection{Proponowane biomarkery dojrzałości funkcjonalnej}
% %,----------------------------------------------------------

% Wybór biomarkerów do~modelu oceny dojrzałości funkcjonalnej
% opiera się na~czterech kryteriach: (1)~związku
% z~konkretnym mechanizmem biologicznym dojrzewania miodu,
% (2)~zwalidowaniu analitycznym lub~zaawansowanym stadium
% walidacji, (3)~dostępności technicznej i~ekonomicznej
% w~laboratoriach kontrolnych, (4)~niezależności od~warunków
% klimatycznych i~systemu produkcji. Tabela~\ref{tab:biomarkers}
% przedstawia zestawienie proponowanych biomarkerów wraz
% z~ich charakterystyką.

% \begin{table}[H]
% % \begin{adjustwidth}{-\extralength}{0cm}
% \centering
% \caption{Proponowane biomarkery dojrzałości funkcjonalnej miodu w~modelu MODF.
%          Źródła: opracowanie własne na~podstawie~\cite{ruegg1981,zhang2021,sun2021,nmrhoney2023,invertase2020,schoder2026}.}
% \label{tab:biomarkers}
% \small
% \begin{tabular}{p{3.2cm}p{2.2cm}p{2.8cm}p{3.4cm}p{1.8cm}}
% %\begin{tabularx}{\fulllength}{LCCCC}
% \hline
% \textbf{Biomarker} &
% \textbf{Mechanizm biologiczny} &
% \textbf{Metoda oznaczania} &
% \textbf{Wartość progowa / zakres} &
% \textbf{Status walidacji} \\
% \hline
% Aktywność wody ($a_w$)
%     & Mikrobiologiczna stabilność osmotyczna
%     & Elektryczna metoda pojemnościowa (Rotronic, AquaLab)
%     & $a_w < 0{,}60$ (dojrzały); $> 0{,}62$ (ryzyko fermentacji)
%     & Zwalidowana \\
% \hline
% Wilgotność (\%H$_2$O)
%     & Odparowanie wody z~nektaru
%     & Refraktometria (n. Abbego)
%     & $\leq 20{,}0\%$ (Codex); $\leq 18{,}0\%$ (zalecane)
%     & Zwalidowana \\
% \hline
% Aktywność diastazy (DN)
%     & Aktywność amylaz ze~ślinianek pszczoły
%     & Phadebas, spektrofotometria
%     & $\geq 8$~j.~Schade (Codex); $\geq 3$~j. dla~miodów o~niskiej aktywności enzym.
%     & Zwalidowana \\
% \hline
% Aktywność inwertazy ($\alpha$-gluko\-zydazy)
%     & Hydroliza sacharozy; wskaźnik freshness
%     & Spektrofotometria UV/VIS, p-NPG
%     & $\geq 4$~j.~Siegenthaler; obniżona: wskaźnik obróbki termicznej lub~niedojrzałości
%     & Zwalidowana~\cite{invertase2020} \\
% \hline
% Aktywność oksydazy glukozowej
%     & Biosynteza H$_2$O$_2$ i~kwasu glukonowego
%     & Spektrofotometria; test peroksydazowy
%     & Wartości referencyjne: 0--11,2~$\mu$g H$_2$O$_2$/g$\cdot$h (raw honey)~\cite{invertase2020}; max.\ przed pełnym zasklepieniem
%     & Częściowo zwalidowana \\
% \hline
% Profil kwasów organicznych (glukonian, mleczan, cytrynian, jabłczan)
%     & Metabolizm enzymatyczny i~bakteryjny; wskaźnik dojrzałości metabolomicznej
%     & HPLC-DAD, GC-MS
%     & Glukonian: 3,72--8,29~g/kg (zal. od~typu botan.)~\cite{nature2025}; mleczan: marker aktywności LAB w~miodach bezżądłowych
%     & Zwalidowana (HPLC-DAD)~\cite{orgacids2025} \\
% \hline
% Stosunek fruktozy do~glukozy (F/G)
%     & Profil cukrowy jako wynik hydrolizy sacharozy
%     & NMR, HPLC
%     & $\geq 0{,}9$ dla~miodu dojrzałego; $< 0{,}8$ wskazuje na~niedojrzałość lub~adulterację
%     & Zwalidowana \\
% \hline
% \textsuperscript{1}H-NMR fingerprinting
%     & Kompleksowy obraz składu chemicznego (cukry, aminokwasy, kwasy, polifenole)
%     & \textsuperscript{1}H-NMR 400--600~MHz; bazy: Bruker HMDB, QC-Expert
%     & Klasyfikacja botaniczna, geograficzna, wykrywanie adulteracji syropowej
%     & Zwalidowana~\cite{nmrhoney2023,schoder2026} \\
% \hline
% Kwas dekenodwukarboksylowy (\textit{decenedioic acid})
%     & Kwas tłuszczowy pszczelego pochodzenia; gromadzony w~miodzie dojrzałym
%     & GC-MS, LC-MS
%     & Istotnie wyższe stężenie w~miodzie dojrzałym vs.~niedojrzałym ($p < 0{,}01$)~\cite{sun2021}; 3 typy botaniczne (rzepak, akacja, jujube)
%     & Częściowo zwalidowana~\cite{sun2021} \\
% \hline
% Liczba drożdży (CFU/g)
%     & Potencjał fermentacyjny (zależny od~$a_w$, $T$, $t$)
%     & Metoda płytkowa (DRBC Agar, ISO 21527-1)
%     & Brak limitu normatywnego; $> 1000$~CFU/g + $a_w > 0{,}62$: ryzyko wysokie
%     & Zwalidowana metodycznie, lecz niewystarczająca jako izolowany parametr \\
% \hline
% %\end{tabularx}
% \end{tabular}
% % \end{adjustwidth}
% \end{table}

\subsubsection{Aktywność wody jako parametr priorytetowy}

Spośród wszystkich wymienionych biomarkerów aktywność wody
($a_w$) jest proponowana jako parametr priorytetowy z~trzech
powodów. Po~pierwsze, jest ona thermodynamiczną miarą
\emph{dostępności wody dla~drobnoustrojów}, a~nie jedynie
zawartości wody w~masie produktu; bezpośrednio koreluje
z~biologiczną możliwością wzrostu drożdży osmofilnych,
co~czyni ją~jedynym parametrem, który może zastąpić
całą gamę pośrednich wskaźników mikrobiologicznych~\cite{ruegg1981,waaw2006}.
Po~drugie, pomiar $a_w$ jest szybki (3-5~minut), tani
i~możliwy do~wykonania w~pasiece czy laboratorium kontrolnym
standardowym przyrządem typu AquaLab CX3TE lub~Rotronic
HygroLab; wynik jest powtarzalny i~jednoznaczny,
nie~wymaga interpretacji eksperckiej~\cite{zamora2006}.
Po~trzecie, kryterium $a_w < 0{,}60$ jest oparte
na~biologicznie uzasadnionym mechanizmie, jest progiem
fizjologicznym drożdży osmofilnych wyznaczonym
in~vitro w~warunkach izobarycznych~\cite{waaw2006,ruegg1981}
, a~nie na~arbitralnie wybranej wartości administracyjnej.

Należy przeprowadzić jednak badania czy podobnie jak dla dopuszczalnych wyjątków w zakresie zawartości wody i w tym względzie nie należy ustanowić wyjątków. Stabilność mikrobiologiczna pewnych typów miodów może wynikać nie tyle z niskiej aktywności wody, a ze specyficznych składników obecnych w nektarze roślin. 

% \subsubsection{Aktywność inwertazy i oksydazy glukozowej jako wskaźniki freshness i dojrzałości}

% Aktywność inwertazy ($\alpha$-glukozydazy) jest w~aktualnej
% literaturze uznawana za~czulszy wskaźnik jakości miodu
% niż aktywność diastazy, gdyż inwertaza ulega inaktywacji
% termicznej przy niższej temperaturze i~krócej
% niż~diastaza, przez co~jej obniżona aktywność może
% sygnalizować zarówno przetwarzanie termiczne produktu,
% jak i~jego niedojrzałość~\cite{invertase2020,kostic2022}.
% Badanie Chen i~in.~\cite{zhang2021} wykazało, że~aktywność
% inwertazy wzrasta gradientowo od~0 do~wartości stabilnych
% w~trakcie naturalnego dojrzewania, przy czym tempo wzrostu
% jest najwyższe w~pierwszych 4--8 dniach procesu. Aktywność
% oksydazy glukozowej wykazuje odmienny, nielinearny profil
% kinetyczny: osiąga maksimum tuż przed zasklepieniem
% komórek, a~następnie nieznacznie spada w~miodzie
% zasklepionym~\cite{zhang2021,invertase2020}, co~potwierdza,
% że~zasklepienie nie~jest momentem kulminacyjnym aktywności
% enzymatycznej, lecz jednym z~etapów ciągłego procesu
% biologicznego (por.~Sekcja~3).

% \subsubsection{Profil kwasów organicznych, walidacja HPLC-DAD}

% Niedawna publikacja z~2025~roku~\cite{orgacids2025} opisuje
% w~pełni zwalidowaną metodę HPLC-DAD dla~oznaczania
% 12~kwasów organicznych w~miodzie, osiągając separację
% wszystkich analitów w~ciągu 15~minut, z~LOD 0,01--4,45~$\mu$g/ml
% i~LOQ 0,03--13,50~$\mu$g/ml. Dominującym kwasem organicznym
% we~wszystkich typach miodu jest kwas glukonowy (manuka:
% 3,72~g/kg, lucerna: 8,29~g/kg)~\cite{nature2025},
% podczas gdy kwas mlekowy stanowi marker aktywności
% bakterii kwasu mlekowego (LAB), istotny zwłaszcza
% dla~miodów bezżądłowych. W~kontekście modelu MODF,
% profil kwasów organicznych dostarcza wielowymiarowego
% wektora cech metabolomicznych, który pozwala różnicować
% typy botaniczne i~etapy dojrzewania niezależnie
% od~systemu produkcji.

% \subsubsection{\texorpdfstring{\textsuperscript{1}H-NMR fingerprinting}{1H-NMR fingerprinting}, złoty standard autentyczności}

% \textsuperscript{1}H-NMR w~modelu MODF pełni rolę
% narzędzia weryfikacyjnego na~poziomie III (patrz niżej),
% lecz ze~względu na~wyjątkową wartość diagnostyczną zasługuje
% na~szczegółowe omówienie. NMR pozwala na~jednoczesną
% kwantyfikację cukrów (fruktozy, glukozy, sacharozy,
% erlozy, turanoza), kwasów organicznych, aminokwasów,
% polifenoli i~związków lotnych w~jednym eksperymencie,
% osiągając klasyfikację botaniczną i~geograficzną miodów
% z~dokładnością powyżej 95\% dla~miodów z~dobrze
% reprezentowanymi bazami referencyjnymi~\cite{nmrhoney2023,schoder2026}.
% Platforma Bruker FoodScreener i~analogiczne systemy komercyjne
% (Eurofins FoodAuthentics, QC-Expert MBiotech) są aktualnie
% stosowane przez kilkanaście europejskich laboratoriów
% referencyjnych jako metoda przesiewowa kontroli
% granicznych~\cite{schoder2026}.

% Badanie Tsagkaris i~in.~\cite{tsagkaris2021} podkreśla jednak,
% że~\textsuperscript{1}H-NMR ma~ograniczoną moc różnicującą
% dla~miodów z~regionów niereprezentowanych w~bazach
% referencyjnych. Przegląd z~2025~roku~\cite{schoder2026}
% formułuje to~wprost: \textit{,,high analytical resolution alone
% therefore does not guarantee unambiguous fraud attribution''}
% , co~oznacza, że~wynik NMR dla~miodu etiopskiego z~ula
% kłodowego będzie niepewny bez~etiopskiej bazy referencyjnej.
% Wniosek ten wzmacnia postulat z~Sekcji~6 dotyczący rozszerzenia
% baz EURL o~próbki z~systemów tropikalnych i~tradycyjnych.

%,----------------------------------------------------------
\subsection{Hierarchiczny model decyzyjny}
%,----------------------------------------------------------

Na~podstawie opisanych biomarkerów proponujemy hierarchiczny
model decyzyjny (HMD) składający się z~trzech poziomów,
zorganizowanych zgodnie z~zasadą proporcjonalności
, każdy kolejny poziom jest bardziej kosztowny analitycznie
i~stosowany wyłącznie wtedy, gdy poprzedni nie~daje
jednoznacznego wyniku lub~gdy wynik Poziomu~I jest
pozytywny (tj.~kryteria są spełnione):




% \begin{figure}[H]
% \centering % Wyśrodkowanie rysunku
% \begin{tikzpicture}[
%   node distance=2.2cm, % Automatyczne odstępy między węzłami
%   level/.style = {
%     draw, 
%     rectangle, 
%     rounded corners, 
%     minimum width=9cm, 
%     text width=8.5cm,      
%     minimum height=1.2cm, 
%     align=center, % Lepiej zastępuje "text centered" przy łamaniu linii
%     font=\small
%   },
%   arrow/.style = {->, thick, >=stealth}
% ]
%   % Definicja węzłów
%   \node[level, fill=green!15]  (L1) 
%     {\textbf{Poziom I (przesiewowy):} Wilgotność $\leq 20\%$ \textsc{oraz} $a_w < 0{,}60$ \\ 
%      $\Rightarrow$ Brak dyskwalifikacji niezależnie od~CFU/g};

%   \node[level, fill=yellow!20, below of=L1] (L2) 
%     {\textbf{Poziom II (jakościowy):} Diastaza $\geq 8$~j.~Sch., Inwertaza $\geq 4$~j., F/G $\geq 0{,}9$, HMF $\leq 40$~mg/kg \\ 
%      $\Rightarrow$ Ocena dojrzałości funkcjonalnej};

%   \node[level, fill=orange!20, below of=L2, node distance=2.5cm] (L3) 
%     {\textbf{Poziom III (arbitrażowy):} NMR fingerprinting + metabolomika + analiza łańcucha dostaw \\ 
%      $\Rightarrow$ Rozróżnienie: stabilizacja vs.\ fałszowanie};

%   % Strzałki - naprawiona składnia (-- zamiast przecinka)
%   \draw[arrow] (L1) -- (L2) node[midway, right, xshift=0.1cm, text width=4cm, font=\scriptsize, align=left] 
%     {wynik graniczny lub\\CFU/g $> 1000$};
    
%   \draw[arrow] (L2) -- (L3) node[midway, right, xshift=0.1cm, text width=4cm, font=\scriptsize, align=left] 
%     {wyniki rozbieżne lub\\podejrzenie adulteracji};
% \end{tikzpicture}
% \caption{Schemat wielopoziomowej oceny jakości i autentyczności miodu.} % Dobra praktyka: dodanie podpisu
% \label{fig:hmd}
% \end{figure}


\subsubsection{Poziom I – Podstawowy (przesiewowy)}

Poziom I stanowi filtr pierwszego rzędu oparty na~dwóch
parametrach mierzonych bezpośrednio na~produkcie:

\begin{itemize}
    \item \textbf{Wilgotność $\leq 20\%$} przy uwzględnieniu prawnie określonych wyjątków (refraktometria,
          metoda Codex STAN 12-1981), parametr normatywny,
          obligatoryjny we~wszystkich głównych standardach;
    \item \textbf{Aktywność wody $a_w < 0{,}60$} przy uwzględnieniu możliwych wyjątków, parametr
          proponowany jako uzupełnienie normatywne; przy spełnieniu
          obu tych kryteriów miód jest \emph{mikrobiologicznie
          stabilny z~mocy stanu wiedzy mikrobiologii}, wzrost drożdży
          jest niemożliwy niezależnie od~ich liczebności~\cite{ruegg1981,waaw2006}.
\end{itemize}

\textbf{Zasada kluczowa poziomu I:} jeżeli wilgotność $\leq 20\%$
\textsc{ORAZ} $a_w < 0{,}60$, produkt nie~podlega dyskwalifikacji
na~podstawie liczby drożdży (CFU/g), niezależnie od~stwierdzonej
wartości tego parametru. Dyskwalifikacja na~podstawie samego
CFU/g, bez~potwierdzenia przekroczenia progu $a_w$, jest
nieuzasadniona naukowo i~nieadekwatna prawnie~\cite{codex_stan12,eudir2001110}.

Jeżeli Poziom I wskazuje na~stan graniczny ($18{,}1$--$20{,}0\%$
wilgotności lub~$a_w = 0{,}60$--$0{,}62$) lub~liczba drożdży
przekracza $1000$~CFU/g, analiza jest eskalowana do~Poziomu~II.

\subsubsection{Poziom II – Rozszerzony (jakościowy)}

Poziom II dostarcza oceny dojrzałości funkcjonalnej
we~wszystkich trzech wymiarach zdefiniowanych w~Sekcji~3.3:
stabilności mikrobiologicznej, chemicznej i~biologicznej.
Parametry Poziomu~II obejmują:

\begin{enumerate}
    \item \textbf{Profil enzymatyczny:}
    \begin{itemize}
        \item Aktywność diastazy $\geq 8$~j.~Schade (lub $\geq 3$~j.
              dla~typów botanicznych z~naturalnie niską aktywnością
              amylolityczną, np.~akacja, cytrusy, lawenda)~\cite{codex_stan12};
      %   \item Aktywność inwertazy $\geq 4$~j.~Siegenthaler;
      %         wartości poniżej 4~j. wskazują na~niedojrzałość
      %         lub~przetwarzanie termiczne, wymaga korelacji
      %         z~HMF~\cite{invertase2020};
      %   \item Aktywność oksydazy glukozowej, marker biologiczny
      %         aktywności antyoksydacyjnej i~antybakteryjnej;
      %         wartości referencyjne dla~danego typu botanicznego~\cite{invertase2020};
    \end{itemize}
    \item \textbf{Parametry chemiczne:}
    \begin{itemize}
        \item Stosunek F/G $\geq 0{,}9$; HMF $\leq 40$~mg/kg;
              kwasowość wolna $\leq 50$~mEq/kg (Codex)~\cite{codex_stan12};
      %   \item Profil kwasów organicznych (opcjonalnie):
      %         kwas glukonowy jako dominujący; relacja
      %         mleczan/glukonian jako marker aktywności LAB~\cite{orgacids2025};
    \end{itemize}
%     \item \textbf{Parametry biologiczne:}
%     \begin{itemize}
%         \item Aktywność antyoksydacyjna (DPPH, FRAP)
%               w~porównaniu z~wartościami referencyjnymi
%               dla~danego typu botanicznego~\cite{guo2020};
%         \item Zawartość polifenoli ogółem (metoda Folin-Ciocalteu)
%               jako wskaźnik integralności biologicznej produktu~\cite{guo2020}.
%     \end{itemize}
\end{enumerate}

% Jeżeli Poziom~II dostarcza wyników rozbieżnych (np.~niska
% aktywność diastazy przy niskim HMF i~prawidłowej aktywności
% inwertazy, co~może wskazywać na~typ botaniczny o~niskiej
% aktywności amylolitycznej) lub~jeśli wyniki sugerują
% potencjalne fałszowanie, analiza przechodzi do~Poziomu~III.

% \subsubsection{Poziom III – Arbitrażowy}

% Poziom III jest przeznaczony wyłącznie do~rozstrzygania
% przypadków spornych i~weryfikacji podejrzenia celowego
% fałszowania. Jest on~kosztowny analitycznie i~wymaga
% dostępu do~specjalistycznej infrastruktury laboratoryjnej;
% nie~jest wskazany jako rutynowy element kontroli
% jakości w~punkcie odbioru. Narzędzia Poziomu~III obejmują:

% \begin{enumerate}
%     \item \textbf{\textsuperscript{1}H-NMR fingerprinting}
%           (400--600~MHz), kompleksowa analiza składu,
%           klasyfikacja botaniczna i~geograficzna,
%           wykrywanie adulteracji syropowej; wynik porównywany
%           z~bazą referencyjną EURL lub~Bruker FoodScreener~\cite{nmrhoney2023,schoder2026};
%     \item \textbf{Izotopowa analiza masowa (IRMS):}
%           $\delta^{13}$C, wykrywanie dodatku syropów C$_4$
%           (kukurydza, trzcina cukrowa); LC-IRMS,
%           izotopowe profilowanie poszczególnych cukrów
%           z~dokładnością do~$\pm 0{,}5$\textperthousand~\cite{nmrhoney2023};
%     \item \textbf{Metabolomika ukierunkowana i~nieukierunkowana}
%           (UPLC-QTOF-MS, GC-MS), identyfikacja biomarkerów
%           dojrzałości molekularnej (kwas dekenodwukarboksylowy,
%           turanoza, lotne estry i~aldehydy)~\cite{sun2021,chromatometab2024};
%     \item \textbf{Analiza łańcucha dostaw (traceability audit)}:
%           weryfikacja dokumentacji wilgotności i~temperatury
%           \textit{post-harvest}; dane logger'ów temperaturowych;
%           certyfikaty analizy z~pasieki; deklaracje pszczelarskie
%           dotyczące zastosowania dehumidyfikacji~\cite{cbi_traceability2024}.
% \end{enumerate}

% Wynik Poziomu~III powinien umożliwiać trójkategoryczną
% klasyfikację produktu: (a)~miód dojrzały funkcjonalnie
% , wolny od~ograniczeń rynkowych; (b)~miód z~zastosowaną
% technologiczną dehumidyfikacją, dopuszczony do~obrotu
% z~deklaracją procesu; (c)~miód adulterowany lub~niedojrzały
% systemowo, dyskwalifikowany~\cite{garcia2025,codex_stan12}.

%,----------------------------------------------------------
\subsection{Rekomendacje dla organów regulacyjnych i Apimondia}
%,----------------------------------------------------------

Na~podstawie proponowanego modelu MODF i~przeprowadzonej
analizy naukowej, niniejszy artykuł formułuje następujące
rekomendacje adresowane do~kluczowych interesariuszy:

\subsubsection{Rekomendacje dla Apimondia}

\begin{enumerate}
    \item \textbf{Uzupełnić stanowisko z~2023~roku o~kryterium
          $a_w$:} rekomendacja Apimondia powinna explicite
          wskazywać, że~obecność drożdży osmofilnych
          w~miodzie nie~stanowi przesłanki do~stwierdzenia
          niedojrzałości, jeśli wilgotność $\leq 20\%$
          i~$a_w < 0{,}60$; oba parametry powinny być
          wymienione jako jednoczesne warunki
          konieczne~\cite{ruegg1981,waaw2006};
%     \item \textbf{Przyjąć wieloparametryczną definicję
%           niedojrzałości:} niedojrzały jest miód, który jednocześnie
%           wykazuje wilgotność $> 20\%$, aktywność wody
%           $a_w > 0{,}62$, obniżony profil enzymatyczny
%           i~brak stabilności chemicznej, nie~zaś miód,
%           który wyłącznie wykazuje liczbę drożdży powyżej
%           arbitralnego progu;
    \item \textbf{Uwzględnić regionalne zróżnicowanie systemów
          produkcji:} stanowisko powinno zawierać
          wyraźną klauzulę dotyczącą miodów z~klimatu
          tropikalnego i~systemów tradycyjnych,
          uznającą dehumidyfikację technologiczną
          za~legalną praktykę produkcyjną~\cite{tadele2025} lub wykluczać konkretne techniki dehumidyfikacji, które negatywnie wpływają na produkt;
    \item \textbf{Zainicjować dialog z~ISO TC~34/SC~17:}
          norma ISO 24607:2025 dotycząca miodu pszczelego
          powinna zostać rozszerzona o~kryterium $a_w$
          jako parametru normatywnego;~\cite{codex_stan12}.
\end{enumerate}

\subsubsection{Rekomendacje dla organów UE (KE, DG SANTE)}

\begin{enumerate}
    \item \textbf{Wdrożyć $a_w$ jako obligatoryjny parametr
          kontrolny} podczas urzędowych kontroli jakości miodu w tym każdorazowo przez inspekcje graniczne dla~partii miodu
          importowanego z~krajów tropikalnych (Afryka Wschodnia,
          Azja Południowo-Wschodnia, Ameryka Łacińska);          
%     \item \textbf{Rozszerzyć projekt European Honey Platform
%           (uruchomiony w~2024~roku) o~globalne bazy
%           referencyjne NMR i~metabolomiczne,} reprezentujące
%           systemy produkcji tropikalnej i~tradycyjnej;
%           bez~tych baz zaawansowane metody omiczne nie~mogą
%           pełnić funkcji narzędzi kontrolnych dla~importu
%           z~tych regionów~\cite{schoder2026,honeyid2025};
%     \item \textbf{Wypracować protokół oceny warunków
%           \emph{post-harvest}:} inspekcje graniczne powinny
%           uwzględniać dokumentację temperaturową i~wilgotnościową
%           łańcucha dostaw; fermentacja stwierdzona
%           w~chwili kontroli nie~powinna być automatycznie
%           interpretowana jako wada produkcyjna bez~analizy
%           warunków transportu i~magazynowania~\cite{analytica2020,apinz2021};
%     \item \textbf{Promować wdrożenie NMR i~LC-IRMS
%           jako metod pierwszorzędowych} w~inspekcjach
%           granicznych dla~partii powyżej 5~ton, zastępując
%           lub~uzupełniając kosztowną metodę SNIF-NMR;
%           zastosowanie platformy Bruker FoodScreener
%           pozwala na~wynik w~ciągu 30~minut i~jest
%           kalibrowalne dla~nowych typów geograficznych~\cite{nmrhoney2023,schoder2026}.
\end{enumerate}

\subsubsection{Rekomendacje dla laboratoriów i środowiska naukowego}

\begin{enumerate}
    \item \textbf{Standaryzować protokół wieloparametrycznej
          oceny mikrobiologicznej} obejmujący jednoczesny
          pomiar CFU/g, $a_w$, temperatury i~stanu fazowego
          miodu (płynny/skrystalizowany); wyniki powinny
          być raportowane łącznie, nie~jako izolowane
          wartości~\cite{zamora2006,analytica2020};
    \item \textbf{Walidować wielolaboratoryjnie biomarker
          kwasu dekenodwukarboksylowego} i~rozszerzyć jego
          weryfikację poza trzy typy botaniczne (rzepak,
          akacja, jujube) na~miody tropikalne, stingless bee
          i~tradycyjne systemy produkcji~\cite{sun2021};
    \item \textbf{Rozwijać przenośne narzędzia analizy \textit{in situ}:}
          miniaturyzowane spektrometry NIR i~Raman
          pozwalają na~pomiar wilgotności, $a_w$
          i~profilu cukrowego bezpośrednio w~pasiece
          lub~na~granicy, z~dokładnością porównywalną
          z~metodami laboratoryjnymi~\cite{honeyid2025,schoder2026}.
\end{enumerate}